\section{Localized five-dimensional sectors and their Higgs coordinates}
\label{sec:local-sectors}

The compact argument requires the holomorphic flavor-moment-map coordinates
that coincide with geometric deformation parameters. We identify these
coordinates as affine schemes; the remaining Higgs-branch data of the
rank-one theories will not enter the global charge map.

\subsection{Geometric engineering and the meaning of the projection}
\label{subsec:local-conventions}

M-theory on a noncompact canonical Calabi--Yau threefold singularity $Y$
defines a five-dimensional theory $\cT_Y$
\cite{Seiberg:1996,Douglas:1996,Intriligator:1997}. A crepant resolution of $Y$
parametrizes its Coulomb data, while normalizable complex deformations encode
Higgs-branch data. This dictionary has been developed in many examples and
can be refined to the stratification by symplectic leaves
\cite{Closset:2020afy,DeMarco:2026higgs}.

Two qualifications will be important. First, the base $\Def(Y)$ of complex
deformations is not itself the hyperk\"ahler Higgs branch. In the examples
below it is obtained from a chosen complex structure on the Higgs cone by
taking affine invariants under a complexified flavor torus. It therefore
records Cartan moment-map values but forgets the remaining orbit coordinates.
Second, strict $L^2$ normalizability in an isolated noncompact metric is a
separate analytic question. The compact theorem uses limits of modes on a
compact threefold, where their norms are finite. We return to noncompact
limits in Section~\ref{sec:decoupling}.

The three local dictionaries are summarized in
Table~\ref{tab:local-dictionary}. Here $\cB_{\mathrm{red}}$ denotes the reduced
geometric deformation component selected by an ordinary complex Higgs vev.

\begin{table}[t]
\centering
\small
\begin{adjustbox}{max width=\textwidth,center}
\begin{tabular}{@{}
>{\raggedright\arraybackslash}p{0.18\textwidth}
>{\raggedright\arraybackslash}p{0.22\textwidth}
>{\raggedright\arraybackslash}p{0.20\textwidth}
>{\raggedright\arraybackslash}p{0.27\textwidth}@{}}
\toprule
singularity & local 5d sector & $\cB_{\mathrm{red}}$ & invariant moment-map coordinates \\
\midrule
ordinary double point & free hypermultiplet & $\CC_t$ & $t=h\widetilde h$ \\
cone over $S_7=\operatorname{Bl}_2\mathbb P^2$ & rank-one $E_2$ theory & $\CC_\beta$ & $\beta=J_0$, the $SU(2)$ Cartan moment map \\
cone over $S_6=\operatorname{Bl}_3\mathbb P^2$ & rank-one $E_3$ theory & $\CC^2_{H_1,H_2}\cup_{0}\CC_L$ & $(H_1,H_2)$ for $SU(3)$ and $L$ for $SU(2)$ \\
\bottomrule
\end{tabular}
\end{adjustbox}
\caption{The local holomorphic projection used in the compact theorem. The
$E_2$ miniversal base also contains an additional nilpotent tangent direction
generated by $U$; it is not an ordinary reduced Higgs trajectory.}
\label{tab:local-dictionary}
\end{table}

\subsection{The conifold sector}
\label{subsec:conifold-sector}

For the ordinary double point
\begin{equation}
 Y_0=\{xy-zw=0\}\subset\CC^4,
 \label{eq:conifold}
\end{equation}
the small resolution replaces the singular point by a curve
$C\cong\mathbb P^1$. An M2-brane wrapping $C$ gives a five-dimensional
hypermultiplet. In one complex structure let $(h,\widetilde h)$ be its two
complex scalars, with opposite charges under the relevant $U(1)$. The local
complex moment map and the conifold deformation coordinate are the same
invariant,
\begin{equation}
 t=h\widetilde h,
 \qquad
 Y_t=\{xy-zw=t\}.
 \label{eq:conifold-map}
\end{equation}
This local map can also be obtained from the type-IIA description of the
rank-zero theory \cite{Collinucci:2021jgb}.

For later use, every $t\in\CC$ has a representative with zero real moment
map. If $t=|t|e^{\ii\vartheta}$, take
\begin{equation}
 h=\sqrt{|t|}e^{\ii\vartheta},
 \qquad
 \widetilde h=\sqrt{|t|}.
 \label{eq:conifold-balanced}
\end{equation}
Then $h\widetilde h=t$ and
$|h|^2-|\widetilde h|^2=0$. This elementary observation is the local reason
that arbitrary complex coefficients in a homology relation can solve the full
triplet of moment-map equations.

\subsection[The E2 reduced component and nilpotent tangent]{The $E_2$ reduced component and nilpotent tangent}
\label{subsec:e2-sector}

The cone over $S_7$ engineers the rank-one $E_2$ fixed point with flavor
algebra
\begin{equation}
 \mathfrak e_2=\mathfrak{su}(2)\oplus\mathfrak u(1).
\end{equation}
In a chosen complex structure, the Higgs chiral ring is generated by an
$SU(2)$ moment-map triplet $(J_+,J_0,J_-)$ and a further current-multiplet
operator $U$, denoted $P$ in \cite{Cremonesi:2015lsa}:
\begin{equation}
 \CC[\cH_{E_2}]
 =\frac{\CC[J_+,J_0,J_-,U]}
 {(J_0^2-J_+J_-,\ U^2,\ UJ_+,\ UJ_0,\ UJ_-)}.
 \label{eq:e2-full-ring}
\end{equation}
The reduction is the $A_1$ surface
\begin{equation}
 \cH_{E_2,\mathrm{red}}
 =\{J_0^2=J_+J_-\}
 \cong\CC^2/\mathbb Z_2,
 \label{eq:e2-reduced}
\end{equation}
and $U$ defines an embedded point at its vertex.

The complexified $SU(2)$ Cartan has weights $(2,0,-2,0)$ on
$(J_+,J_0,J_-,U)$. A weight-zero monomial in the reduced ring is a power
of $J_0$, because
\begin{equation}
 J_+^aJ_0^bJ_-^a=(J_+J_-)^aJ_0^b=J_0^{2a+b}.
\end{equation}
Every nonzero monomial containing $U$ is proportional to $U$. It follows that
\begin{equation}
 \CC[\cH_{E_2}]^{\CC^*}
 =\frac{\CC[J_0,U]}{(U^2,UJ_0)}.
 \label{eq:e2-invariants}
\end{equation}
Altmann's miniversal base for the degree-seven cone is
\begin{equation}
 \Def(C(S_7))
 =\Spec\frac{\CC[s_2,s_1]}{(s_1^2,s_1s_2)}
 \label{eq:e2-altmann}
\end{equation}
\cite{Altmann:1997}. The root marking from the crepant resolution fixes the
dictionary, up to nonzero normalization,
\begin{equation}
 \beta=s_2\longleftrightarrow J_0,
 \qquad
 \alpha=s_1\longleftrightarrow U.
 \label{eq:e2-dictionary}
\end{equation}
Thus the complete nonreduced schemes agree, not only their tangent dimensions.
The reduced line $\alpha=0$ is the physical finite Higgs component. Since
$U^2=UJ_0=0$, the $U$ direction exists over dual numbers but vanishes at every
ordinary complex point; it cannot be counted as a second finite branch.

The Cartan interpretation is explicit in the double-cover coordinates:
\begin{equation}
 J_+=a^2,
 \qquad J_0=ab,
 \qquad J_-=b^2,
 \qquad (a,b)\sim(-a,-b).
 \label{eq:e2-double-cover}
\end{equation}
The Cartan acts as $(a,b)\mapsto(\lambda a,\lambda^{-1}b)$ and has moment
maps
\begin{equation}
 \mu_{E_2}^{\CC}=ab=\beta,
 \qquad
 \mu_{E_2}^{\mathbb R}=\frac12(|a|^2-|b|^2).
 \label{eq:e2-moment-maps}
\end{equation}
The choice $a=\sqrt{|\beta|}e^{\ii\arg\beta}$,
$b=\sqrt{|\beta|}$ realizes any complex $\beta$ while setting the real
moment map to zero.

\subsection[The two components of the E3 sector]{The two components of the $E_3$ sector}
\label{subsec:e3-sector}

The cone over $S_6$ engineers the rank-one $E_3$ fixed point with flavor
algebra
\begin{equation}
 \mathfrak e_3=\mathfrak{su}(3)\oplus\mathfrak{su}(2).
\end{equation}
Its reduced Higgs branch is the union
\begin{equation}
 \cH_{E_3}
 =\underbrace{\overline{\mathcal O_{\min}(\mathfrak{sl}_3)}}_
 {\cH_{\mathfrak{sl}_3}}
 \mathop{\cup}_{\{0\}}
 \underbrace{\overline{\mathcal O_{\min}(\mathfrak{sl}_2)}}_
 {\cH_{\mathfrak{sl}_2}} .
 \label{eq:e3-components}
\end{equation}
The two components have quaternionic dimensions two and one. In terms of the
complex flavor moment maps
$\Phi\in\mathfrak{sl}_3$ and
$\widetilde T\in\mathfrak{sl}_2$, the chiral-ring relations are
\begin{equation}
 \Phi^2=0,
 \qquad
 \tr(\widetilde T^2)=0,
 \qquad
 \Phi^i{}_j\widetilde T^{\dot\alpha\dot\beta}=0
 \label{eq:e3-chiral-relations}
\end{equation}
\cite{Cremonesi:2015lsa}. The mixed relation says scheme-theoretically that
positive-degree functions on different components annihilate one another.
The intersection is one reduced point.

Let $(H_1,H_2)$ be root-basis Cartan coordinates of $\Phi$ and let $L$ be
the Cartan coordinate of $\widetilde T$. The maximal-torus invariant rings of
the two minimal nilpotent orbit closures are respectively
$\CC[H_1,H_2]$ and $\CC[L]$. Taking invariants preserves the fiber product,
because a complex torus is linearly reductive. Hence
\begin{equation}
 \CC[\cH_{E_3}]^{T_{\CC}}
 =\CC[H_1,H_2]\mathop{\times}_{\CC}\CC[L]
 \cong\frac{\CC[H_1,H_2,L]}{(H_1L,H_2L)}.
 \label{eq:e3-invariants}
\end{equation}
We reserve $\cH_{\mathfrak{sl}_3}$ and $\cH_{\mathfrak{sl}_2}$ for the full
orbit closures and denote their invariant-coordinate images by
\begin{equation}
 \cB_{\mathfrak{sl}_3}:=\Spec\CC[H_1,H_2]\cong\CC^2,
 \qquad
 \cB_{\mathfrak{sl}_2}:=\Spec\CC[L]\cong\CC.
 \label{eq:e3-invariant-images}
\end{equation}
The miniversal base of the degree-six cone is
\begin{equation}
 \Def(C(S_6))
 =\Spec\frac{\CC[s_1,s_2,s_3]}{(s_1s_3,s_2s_3)}.
 \label{eq:e3-altmann}
\end{equation}
The compatible root marking gives
\begin{equation}
 (s_1,s_2)\longleftrightarrow(H_1,H_2)
 \quad\text{on }\cB_{\mathfrak{sl}_3},
 \qquad
 s_3\longleftrightarrow L
 \quad\text{on }\cB_{\mathfrak{sl}_2}.
 \label{eq:e3-dictionary}
\end{equation}
The mixed equations in \eqref{eq:e3-invariants} matter in the compact
problem: at one $E_3$ point, an ordinary Higgs trajectory chooses the
$SU(3)$ or $SU(2)$ moment maps, but not both.

With $K_{S_6}^{\perp}$ denoting the orthogonal complement of the canonical
class in $H_2(S_6;\ZZ)$ for the intersection pairing, the corresponding root
decomposition inside the del Pezzo exceptional divisor is
\begin{equation}
 K_{S_6}^{\perp}\cong A_2\oplus A_1,
 \qquad
 R_1=e_1-e_2,
 \quad R_2=e_2-e_3,
 \quad R_0=\ell-e_1-e_2-e_3.
 \label{eq:e3-roots}
\end{equation}
Here $(R_1,R_2)$ span $A_2$ and $R_0$ spans its orthogonal $A_1$. Root-basis
coordinates are used throughout so that a charge functional with values
$(m_1,m_2)$ on the simple roots couples as $m_1H_1+m_2H_2$, without an
unwritten inverse Cartan matrix.

As in the $E_2$ case, arbitrary complex Cartan values admit representatives
with zero real moment maps. One may use the rank-one
Atiyah--Drinfeld--Hitchin--Manin (ADHM) presentation of a
minimal nilpotent orbit closure: write
$z_i=q_i\widetilde q_i$, impose $\sum_i z_i=0$, and choose
$|q_i|=|\widetilde q_i|=\sqrt{|z_i|}$ with phases reproducing $z_i$.
This sets the ADHM and flavor real moment maps to zero pairwise. The
$\mathfrak{sl}_2$ component is the same construction with two pairs.

\subsection{Magnetic-quiver and Hasse-diagram check}
\label{subsec:e3-magnetic}

The $\mathfrak{sl}_3$ and $\mathfrak{sl}_2$ components have the affine-$A_2$
triangle and affine-$A_1$ double-edge magnetic quivers, after removing the
decoupled diagonal $U(1)$
\cite{Ferlito:2016,Cabrera:2018,Closset:2020afy}; the two-cone structure of
the $E_3$ Higgs branch goes back to \cite{Seiberg:1996,Morrison:1996pp}.
Their Coulomb-branch dimensions are $3-1=2$ and $2-1=1$, agreeing with
\eqref{eq:e3-dictionary}. With $x=t^2$, the Abelian monopole formula
\cite{Cremonesi:2013} gives
\begin{align}
 H_{\mathfrak{sl}_2}(x)
 &=\frac{1+x}{(1-x)^2}
   =\sum_{n\geq0}(2n+1)x^n,
 \label{eq:e3-hs-a1}\\
 H_{\mathfrak{sl}_3}(x)
 &=\frac{1+4x+x^2}{(1-x)^4}
   =\sum_{n\geq0}(n+1)^3x^n.
 \label{eq:e3-hs-a2}
\end{align}
Since the components meet in a reduced point,
\begin{equation}
 H_{E_3}(x)=H_{\mathfrak{sl}_3}(x)+H_{\mathfrak{sl}_2}(x)-1.
 \label{eq:e3-hs-union}
\end{equation}
The coefficient of $x$ is $8+3=11$, the dimension of the
$SU(3)\times SU(2)$ moment-map generators. This is independent of the compact
deformation calculation.

The component data relevant below are collected in
Table~\ref{tab:e3-components}. Each minimal nilpotent orbit closure has only
its open orbit and the origin as symplectic leaves. A generic point therefore
has no transverse interacting endpoint; this local statement does not remove
bulk vectors or other compact sectors.

\begin{table}[t]
\centering
\small
\begin{adjustbox}{max width=\textwidth,center}
\begin{tabular}{@{}lllll@{}}
\toprule
component & magnetic quiver & $\dim_{\HH}$ & coordinate & generic continuous stabilizer \\
\midrule
$\cH_{\mathfrak{sl}_3}$ & affine $A_2$ & $2$ & $(H_1,H_2)$ & $U(1)\times SU(2)$ \\
$\cH_{\mathfrak{sl}_2}$ & affine $A_1$ & $1$ & $L$ & $SU(3)$ \\
\bottomrule
\end{tabular}
\end{adjustbox}
\caption{Independent local data for the two $E_3$ Higgs components. The
$SU(2)$ (respectively $SU(3)$) factor complementary to the active
$\mathfrak{sl}_3$ (respectively $\mathfrak{sl}_2$) orbit is unbroken. The
table lists the full continuous generic stabilizer and omits only finite
subgroups.}
\label{tab:e3-components}
\end{table}
